% !TeX program = xelatex
% !TeX encoding = UTF-8
% Вариант 10. Таблицы пересчитаны по data.txt; генерация H2, seed=42.
\documentclass[12pt,a4paper]{article}

\usepackage{fontspec}
\usepackage[russian]{babel}
\IfFontExistsTF{Times New Roman}
  {\setmainfont{Times New Roman}}
  {\setmainfont{FreeSerif}}
\IfFontExistsTF{Arial}
  {\setsansfont{Arial}}
  {\setsansfont{FreeSans}}
\IfFontExistsTF{Courier New}
  {\setmonofont{Courier New}}
  {\setmonofont{FreeMono}}
\usepackage[left=22mm,right=18mm,top=20mm,bottom=20mm]{geometry}
\usepackage{amsmath,amssymb}
\usepackage{graphicx}
\usepackage{array,booktabs,multirow,tabularx,longtable}
\usepackage{enumitem}
\usepackage{setspace}
\usepackage{caption}
\usepackage{fancyvrb}
\usepackage[hidelinks,unicode]{hyperref}
\usepackage{booktabs}
\usepackage{graphicx}

\hypersetup{
  pdftitle={Обработка результатов измерений: статистический анализ числовой последовательности},
  pdfauthor={Марьин Григорий Алексеевич},
  pdfsubject={Лабораторная работа № 1. Вариант 10}
}
\graphicspath{{./}{figures/}}
\setstretch{1.15}
\setlength{\parindent}{1.25cm}
\setlength{\parskip}{4pt}
\setlength{\emergencystretch}{2em}
\setlength{\tabcolsep}{4pt}
\renewcommand{\arraystretch}{1.2}
\captionsetup{font=small,justification=centering,labelsep=endash}
\setlist[enumerate]{leftmargin=8mm,itemsep=3pt,topsep=4pt}
\setcounter{secnumdepth}{0}
\pagestyle{plain}

\newcommand{\stage}[2]{\subsection*{Этап #1. #2}}
\newcommand{\conclusion}[2]{\par\medskip\noindent\textbf{Вывод из #1 этапа:} #2\par}
\newcommand{\reportfigure}[3]{%
  \begin{center}
    \includegraphics[width=#2\linewidth]{#1}
    \captionof*{figure}{#3}
  \end{center}}
\newcommand{\pmv}[1]{\(\pm #1\)}

\begin{document}

\begin{titlepage}
  \centering
  Федеральное государственное автономное образовательное\newline
  учреждение высшего образования\\[3mm]
  \textbf{«Национальный исследовательский университет ИТМО»}\\[5mm]
  Факультет программной инженерии и компьютерной техники\\[3mm]
  Направление подготовки 09.03.04 «Программная инженерия»\\
  Системное и прикладное программное обеспечение

  \vfill
  {\Large\bfseries Отчёт\par}
  \vspace{4mm}
  {\large по лабораторной работе № 1\par}
  \vspace{4mm}
  {\large\bfseries «Обработка результатов измерений:\\
  статистический анализ числовой последовательности»\par}
  \vspace{4mm}
  По моделированию\\[3mm]
  \textbf{Вариант: 10}
  \vfill

  \begin{flushright}
    \begin{minipage}{0.61\textwidth}
      \raggedleft
      \textbf{Выполнил:}\\
      студент 3 курса\\
      Марьин Григорий Алексеевич\\
      Группа: P3312\\[4mm]
      \textbf{Принял:}\\
      Тропченко Андрей Александрович\\[6mm]
    \end{minipage}
  \end{flushright}

  \vfill
  г. Санкт-Петербург, 2026
\end{titlepage}

\end{document}
